\documentclass[12pt,a4paper]{amsart}

\usepackage{amsmath,amssymb,amsthm}
\usepackage{mathtools}
\usepackage[utf8]{inputenc}
\usepackage[T1]{fontenc}
\usepackage{hyperref}
\usepackage{enumitem,booktabs,array}
\usepackage{geometry}
\geometry{margin=2.5cm}

% -------------------------------------------------------
% Theorem-Umgebungen
% -------------------------------------------------------
\newtheorem{theorem}{Satz}[section]
\newtheorem{lemma}[theorem]{Lemma}
\newtheorem{corollary}[theorem]{Korollar}
\newtheorem{proposition}[theorem]{Proposition}
\theoremstyle{definition}
\newtheorem{definition}[theorem]{Definition}
\newtheorem{example}[theorem]{Beispiel}
\newtheorem{remark}[theorem]{Bemerkung}
\newtheorem{conjecture}[theorem]{Vermutung}

% -------------------------------------------------------
% Eigene Befehle
% -------------------------------------------------------
\newcommand{\N}{\mathbb{N}}
\newcommand{\Z}{\mathbb{Z}}
\newcommand{\Q}{\mathbb{Q}}
\newcommand{\R}{\mathbb{R}}
\newcommand{\C}{\mathbb{C}}
\newcommand{\F}{\mathbb{F}}
\newcommand{\Zp}{\mathbb{Z}_p}
\newcommand{\Qp}{\mathbb{Q}_p}
\newcommand{\A}{\mathbb{A}}
\newcommand{\GL}{\mathrm{GL}}
\newcommand{\SL}{\mathrm{SL}}
\newcommand{\GQ}{\mathrm{Gal}(\bar{\Q}/\Q)}
\newcommand{\Gal}{\mathrm{Gal}}
\newcommand{\Frob}{\mathrm{Frob}}
\newcommand{\Art}{\mathrm{Art}}
\DeclareMathOperator{\ord}{ord}
\DeclareMathOperator{\Hom}{Hom}
\DeclareMathOperator{\tr}{Sp}

% -------------------------------------------------------
\begin{document}
% -------------------------------------------------------

\title{Das Langlands-Programm:\\
       $L$-Funktionen, automorphe Darstellungen\\
       und die globale Korrespondenz}

\author{Michael Fuhrmann}
\address{Specialist Mathematics Project, Build 120}
\email{specialist-maths@localhost}
\date{\today}

\begin{abstract}
Das Langlands-Programm, das Robert Langlands 1967 in einem Brief an André Weil
initiierte, schlägt ein umfangreiches Netz von Vermutungen vor, das Zahlentheorie,
automorphe Formen und Darstellungstheorie verbindet.  Im Kern steht die
\emph{Langlands-Korrespondenz}: eine (vermutete) Bijektion zwischen
$n$-dimensionalen Darstellungen der absoluten Galois-Gruppe $\GQ$ und
automorphen Darstellungen von $\GL_n(\mathbb{A}_\Q)$.  Dieses Paper entwickelt
die notwendigen Grundlagen (Galois-Darstellungen, $L$-Funktionen, automorphe
Formen, lokale Langlands-Korrespondenz) und präsentiert die wichtigsten bekannten
Fälle und Anwendungen, einschließlich Wiles' Beweis des Großen Fermatschen Satzes
als herausragendes Beispiel des Programms.
\end{abstract}

\maketitle
\tableofcontents

% ================================================================
\section{Motivation: Von der quadratischen Reziprozität zu Langlands}
% ================================================================

Das Langlands-Programm verallgemeinert eine lange Kette von Reziprozitätsgesetzen:

\begin{itemize}
  \item \textbf{Quadratische Reziprozität} (Gauß 1796): Für verschiedene ungerade Primzahlen $p, q$:
    $\bigl(\frac{p}{q}\bigr)\bigl(\frac{q}{p}\bigr) = (-1)^{\frac{p-1}{2}\cdot\frac{q-1}{2}}$.
  \item \textbf{Klassenkörpertheorie} (Artin 1927): Jede abelsche Erweiterung $L/K$
    entspricht einer offenen Untergruppe der Idèlenklassengruppe $\mathbb{A}_K^\times/K^\times$.
    Die Artin-Abbildung $\Art_{L/K}$ identifiziert Galois-Gruppen mit Quotienten dieser Gruppe.
  \item \textbf{Langlands}: Nicht-abelsche Erweiterungen entsprechen automorphen
    Darstellungen von $\GL_n$ statt $\GL_1$.
\end{itemize}

Die Schlüsselerkenntnis: Die Arithmetik von Zahlkörpern (Galois-Gruppen) und
die Analysis automorpher Formen (harmonische Analysis auf Adèlen-Gruppen)
sind im Grunde dasselbe Thema.

% ================================================================
\section{Galois-Darstellungen}
% ================================================================

\begin{definition}[$\ell$-adische Galois-Darstellung]
\label{def:galois_darst}
Eine \emph{stetige $\ell$-adische Galois-Darstellung} vom Grad $n$ ist ein
stetiger Gruppenhomomorphismus
\[
  \rho : \GQ \longrightarrow \GL_n(\bar{\Q}_\ell)
\]
(oder nach $\GL_n(\Zp)$, $\GL_n(\Qp)$ etc.) wobei $\bar{\Q}_\ell$ die
$\ell$-adische Topologie trägt.  $\rho$ ist \emph{unverzweigt bei $p$},
falls $\rho|_{I_p} = 1$ für die Trägheitsgruppe $I_p \subset \GQ$.
\end{definition}

\begin{example}[Zyklotomischer Charakter]
Der \emph{zyklotomische Charakter} $\chi_\ell : \GQ \to \Z_\ell^\times$ ist definiert durch
$\sigma(\zeta_{\ell^n}) = \zeta_{\ell^n}^{\chi_\ell(\sigma)}$ für alle $\ell^n$-ten
Einheitswurzeln.  Dies ist eine 1-dimensionale Darstellung.
\end{example}

\begin{example}[Tate-Modul einer elliptischen Kurve]
Sei $E/\Q$ eine elliptische Kurve.  Das \emph{$\ell$-adische Tate-Modul} ist
\[
  T_\ell(E) = \varprojlim_n E[\ell^n](\bar\Q) \;\cong\; \Z_\ell^2.
\]
Die Galois-Wirkung auf den $\ell$-Potenz-Torsionspunkten liefert eine
2-dimensionale Darstellung $\rho_{E,\ell} : \GQ \to \GL_2(\Z_\ell)$.
Für eine Primzahl $p \nmid \ell\cdot\Delta(E)$ (gute Reduktion) gilt:
\[
  \tr(\rho_{E,\ell}(\Frob_p)) = a_p(E) = p + 1 - |E(\F_p)|.
\]
\end{example}

\begin{definition}[$L$-Funktion einer Galois-Darstellung]
\label{def:galois_l}
Für $\rho : \GQ \to \GL_n(\Q_\ell)$, unverzweigt bei fast allen Primzahlen, ist
die \emph{$L$-Funktion}
\[
  L(s, \rho) = \prod_p \det\!\bigl(1 - \rho(\Frob_p)\,p^{-s}\bigr)^{-1}.
\]
\end{definition}

% ================================================================
\section{Automorphe Darstellungen}
% ================================================================

\begin{definition}[Adèlen-Ring]
\label{def:adelen}
Der \emph{Adèlen-Ring} von $\Q$ ist das eingeschränkte Produkt
\[
  \mathbb{A} = \mathbb{A}_\Q = \R \times \prod_p{}' \Q_p
  = \bigl\{(x_\infty, x_2, x_3, \ldots) : x_p \in \Z_p \text{ für fast alle } p\bigr\}.
\]
Es ist ein lokal-kompakter Ring.  Die diagonale Einbettung $\Q \hookrightarrow \mathbb{A}$
macht $\Q$ zu einem diskreten kocompakten Teilring.
\end{definition}

\begin{definition}[Automorphe Darstellung]
\label{def:automorphe_darst}
Eine \emph{automorphe Darstellung} von $\GL_n(\mathbb{A})$ ist eine irreduzible
zulässige Darstellung $\pi = \otimes'_v \pi_v$ (eingeschränktes Tensorprodukt über
alle Stellen $v$), die im Raum $L^2(\GL_n(\Q) \backslash \GL_n(\mathbb{A}), \omega)$
für einen zentralen Charakter $\omega$ vorkommt.  Eine \emph{kuspformige} automorphe
Darstellung hat verschwindende konstante Terme entlang aller echten parabolischen
Untergruppen.
\end{definition}

\begin{example}[Klassische Modulformen als automorphe Darstellungen]
Eine normierte Hecke-Eigenform $f \in S_k(\Gamma_0(N), \chi)$ entspricht
einer kuspformen automorphen Darstellung $\pi_f$ von $\GL_2(\mathbb{A})$.
Die $L$-Funktion von $\pi_f$ ist $L(s, f) = \sum_{n \ge 1} a_n n^{-s}$.
\end{example}

% ================================================================
\section{Die lokale Langlands-Korrespondenz}
% ================================================================

\begin{theorem}[Lokale Langlands-Korrespondenz für $\GL_n$]
\label{thm:lokale_langlands}
Sei $F$ eine endliche Erweiterung von $\Q_p$.  Es gibt eine kanonische Bijektion
\[
  \mathrm{LLC}_F :
  \left\{\parbox{5.5cm}{\centering $n$-dim.\ irreduzible Darstellungen
  der Weil-Deligne-Gruppe $\mathrm{WD}(F)$}\right\}
  \;\longleftrightarrow\;
  \left\{\parbox{5.5cm}{\centering irreduzible zulässige Darstellungen von $\GL_n(F)$}\right\}
\]
mit Kompatibilität der $L$- und $\varepsilon$-Faktoren.
\end{theorem}

\begin{proof}[Status]
Für $\GL_1(F)$: Klassenkörpertheorie. \quad
Für $\GL_2(F)$: Kutzko (1980), Henniart (1993). \quad
Für $\GL_n(F)$ allgemein: Harris--Taylor (2001), Henniart (2000).
\end{proof}

% ================================================================
\section{Die globale Langlands-Korrespondenz}
% ================================================================

\begin{conjecture}[Globale Langlands-Korrespondenz für $\GL_n$]
\label{conj:globale_langlands}
Es gibt eine kanonische Bijektion (\emph{Langlands-Korrespondenz})
\[
  \mathcal{L} :
  \left\{\parbox{5.5cm}{\centering irreduzible $n$-dim.\ $\ell$-adische Darstellungen $\rho$
  von $\GQ$ (halbeinfach, geometrisch)}\right\}
  \;\longleftrightarrow\;
  \left\{\parbox{5.5cm}{\centering kuspiforme automorphe Darstellungen $\pi$
  von $\GL_n(\mathbb{A})$}\right\}
\]
mit $L(s, \rho) = L(s, \pi)$.
\end{conjecture}

\begin{theorem}[Langlands--Tunnell: lösbares Bild]
\label{thm:langlands_tunnell}
Sei $\rho : \GQ \to \GL_2(\C)$ stetig mit lösbarem Bild.  Dann ist $L(s, \rho)$
die $L$-Funktion einer automorphen Form auf $\GL_2(\A)$.
\end{theorem}

% ================================================================
\section{Shimura--Taniyama--Wiles: Ein expliziter Fall}
% ================================================================

\begin{theorem}[Shimura--Taniyama--Wiles--BCDT]
\label{thm:stw}
Jede elliptische Kurve $E/\Q$ ist \emph{modular}: Es gibt eine Hecke-Eigenform
$f_E \in S_2(\Gamma_0(N))$ (für den Führer $N$ von $E$) mit $L(s, E) = L(s, f_E)$.
\end{theorem}

\begin{proof}[Beweisidee (Wiles 1995)]
\begin{enumerate}[label=\arabic*.]
  \item \emph{Restdarstellung}: $\bar\rho_{E,3} : \GQ \to \GL_2(\F_3)$ ist
    modular nach Langlands--Tunnell (da $\GL_2(\F_3)$ lösbar).
  \item \emph{Deformationstheorie}: Der universelle Deformationsring $R$
    und die Hecke-Algebra $T$ sind isomorph ($R = T$, Wiles' Hauptsatz).
  \item \emph{Schluss}: Jede Deformation von $\bar\rho$ (also auch $\rho_{E,3}$
    selbst) kommt von einer Modulform.
\end{enumerate}
Wiles bewies den halbstabilen Fall; Taylor--Wiles ergänzten;
BCDT (2001) vollendeten den allgemeinen Fall.
\end{proof}

\begin{corollary}[Großer Fermatscher Satz (Wiles 1995)]
Die Gleichung $x^n + y^n = z^n$ hat für $n \ge 3$ keine Lösung in positiven
ganzen Zahlen.
\end{corollary}

\begin{proof}[Beweisidee]
Ribet (1990): Eine hypothetische Lösung $(a,b,c)$ würde eine nicht-modulare
elliptische Kurve (Frey-Kurve) liefern — Widerspruch zum Modularity-Theorem.
\end{proof}

% ================================================================
\section{Fontaine-Mazur-Vermutung}
% ================================================================

\begin{conjecture}[Fontaine--Mazur, 1995]
\label{conj:fontaine_mazur}
Eine irreduzible $n$-dimensionale $p$-adische Darstellung $\rho : \GQ \to \GL_n(\Q_p)$
stammt aus der algebraischen Geometrie (d.\,h.\ ist Unterdarstellung der étalen
Kohomologie einer glatten projektiven Varietät über $\Q$) genau dann, wenn sie
\begin{enumerate}[label=(\roman*)]
  \item \emph{geometrisch} ist (unverzweigt außerhalb endlich vieler Primzahlen und
    de Rham bei $p$ im Sinne der $p$-adischen Hodge-Theorie),
  \item \emph{ungerade} ist ($n=2$: komplexe Konjugation hat Eigenwerte $\pm 1$).
\end{enumerate}
Solche geometrischen Darstellungen sollen kuspformen automorphen Darstellungen
entsprechen.
\end{conjecture}

% ================================================================
\section{Bekannte Fälle und offene Probleme}
% ================================================================

\begin{remark}[Übersicht bekannter Fälle]
\begin{center}
\begin{tabular}{l|l|l}
\textbf{Fall} & \textbf{Gruppen} & \textbf{Status} \\
\hline
$n = 1$ (abelsch) & $\GL_1$ & Bewiesen (Klassenkörpertheorie) \\
$n = 2$, lösbares Bild & $\GL_2$ & Bewiesen (Langlands--Tunnell) \\
Elliptische Kurven über $\Q$ & $\GL_2$ & Bewiesen (Wiles + BCDT) \\
Lokale LLC für $\GL_n$ & $\GL_n(F)$, $F/\Q_p$ & Bewiesen (Harris--Taylor, Henniart) \\
Basiswechsel $\GL_n$ & $\GL_n$ & Bewiesen (Arthur--Clozel) \\
$\text{Sym}^m$ für $m \le 4$ & $\GL_2 \to \GL_{m+1}$ & Bewiesen (Kim--Shahidi) \\
$p$-adische LLC für $\GL_2(\Q_p)$ & $\GL_2$ & Bewiesen (Colmez) \\
Geometrische Langlands (Char.\,0) & Kurven über $\C$ & Bewiesen (Gaitsgory et al. 2024) \\
\hline
Globale LLC für $\GL_n$ ($n \ge 2$) & $\GL_n$ & \textbf{OFFEN} \\
Fontaine--Mazur (allgemein) & $\GL_n$ & Teilweise bewiesen \\
Nicht-abelsche Artin-Vermutung & allg.\ $G$ & \textbf{OFFEN} \\
\end{tabular}
\end{center}
\end{remark}

\begin{remark}[Geometrische Langlands 2024]
Im Jahr 2024 vollendete ein Team um Gaitsgory den Beweis der
\emph{geometrischen Langlands-Vermutung} für Kurven über algebraisch
abgeschlossenen Körpern der Charakteristik 0 (de Rham-Setting).
Dies ist ein großer Fortschritt, wenngleich der arithmetische Fall
(Zahlkörper) offen bleibt.
\end{remark}

% ================================================================
\section{Das große Bild}
% ================================================================

\begin{remark}[Wörterbuch]
Das Langlands-Programm schafft ein Wörterbuch:
\begin{center}
\begin{tabular}{l|l}
\textbf{Zahlentheorie / Geometrie} & \textbf{Automorphe Formen / Darstellungstheorie} \\
\hline
Galois-Gruppe $\GQ$ & Algebraische Gruppe $G$ über $\A$ \\
$n$-dim.\ Galois-Darstellung & Automorphe Darstellung von $\GL_n(\A)$ \\
Hasse-Weil-$L$-Funktion & Automorphe $L$-Funktion \\
Verzweigung bei $p$ & Lokale Darstellung $\pi_p$ \\
Frobenius bei $p$ & Hecke-Eigenwert $\lambda_p$ \\
Elliptische Kurve & Gewicht-2-Spitzenform \\
\end{tabular}
\end{center}
Die Vereinigung dieser beiden Welten — Arithmetik und Analysis — ist
das tiefste Thema der modernen Mathematik.
\end{remark}

\begin{thebibliography}{99}

\bibitem{Langlands1970}
R.~P.~Langlands,
\textit{Problems in the theory of automorphic forms},
in: Lectures in Modern Analysis and Applications III,
Lecture Notes in Math.\ \textbf{170}, Springer, 1970, 18--61.

\bibitem{HarrisTaylor2001}
M.~Harris und R.~Taylor,
\textit{The geometry and cohomology of some simple Shimura varieties},
Ann.\ of Math.\ Studies \textbf{151}, Princeton Univ.\ Press, 2001.

\bibitem{Wiles1995}
A.~Wiles,
\textit{Modular elliptic curves and Fermat's Last Theorem},
Ann.\ of Math.\ (2) \textbf{141} (1995), 443--551.

\bibitem{BCDT2001}
C.~Breuil, B.~Conrad, F.~Diamond, R.~Taylor,
\textit{On the modularity of elliptic curves over $\Q$: wild 3-adic exercises},
J.\ Amer.\ Math.\ Soc.\ \textbf{14} (2001), 843--939.

\bibitem{Ribet1990}
K.~Ribet,
\textit{On modular representations of $\Gal(\bar\Q/\Q)$ arising from modular forms},
Invent.\ Math.\ \textbf{100} (1990), 431--476.

\bibitem{ArthurClozel1989}
J.~Arthur und L.~Clozel,
\textit{Simple Algebras, Base Change, and the Advanced Theory of the Trace Formula},
Ann.\ of Math.\ Studies \textbf{120}, Princeton Univ.\ Press, 1989.

\bibitem{Colmez2010}
P.~Colmez,
\textit{Repr\'esentations de $\GL_2(\Q_p)$ et $(\varphi,\Gamma)$-modules},
Ast\'erisque \textbf{330} (2010), 281--509.

\end{thebibliography}

\end{document}
